\documentclass{article}
\usepackage{amsmath,amssymb,hyperref}
\title{What an old custody proof leaves undecided when a new handoff becomes lawful}
\begin{document}
\maketitle

The same two fixed-law proofs can leave different answers after a partner's transfer opens a new timing rule. For a serial custody protocol, the old proof determines an exact interval for each carried expiry: the oldest it might be and the latest it could be. These intervals give exact possible and guaranteed continuations over the corresponding abstract histories. A gate that checks two expiries in one indivisible decision need not have a single earliest compatible expiry vector.

\paragraph{Interface.}
Fix one finite horizon of integer slots \(0,\ldots,H-1\). There are \(k\) named outstanding tags or permissions. The unknown input is the vector \(e\in\{0,\ldots,H\}^k\) of their first usable slots, clipped at either horizon end. The core is a prefix-closed set of labeled timed words accepted by all the other guards: authority transitions, phase, competition at the port, physical-event prerequisites and checks on newly minted receipts. Each step checks zero or one old expiry through \(e_i\le t\); the typed coordinate and slot are determined by the word, not by hidden \(e\). Later guards may depend arbitrarily on earlier labeled steps. A calendar is a special case; an actual partner event can enable a change. This describes feasibility of the labeled words. It does not give the acting crew an observation of a hidden label, an online copy of the true expiry vector, or a distribution on words.

\paragraph{Exact description of an old proof's uncertainty.}
For the admitted language \(S\), define \(u_i\) as the earliest time of a test of \(i\) on any word in \(S\), or \(H\) if absent. Define \(\ell_i\) as one plus the latest time at which a core-legal one-step extension of a word in \(S\) tests \(i\) and is refused, or zero if none. Then the expiry vectors inducing exactly \(S\) are
\[
 \{e : \ell_i\le e_i\le u_i\text{ for every }i\}.
\]
Every accepted test supplies an upper bound. A first refused extension has just one active old-expiry test, giving the indicated lower bound; everything else in that extension passed. Conversely, with those bounds every original admitted word still passes. A new word would have a first previously refused extension, but its lower bound forbids that extension. The box is therefore exact. In particular \(u\) induces precisely \(S\), and two distinct full languages have different upper vectors. The coarsest exact equivalence of expiry histories for this declared interface can thus be represented by \(k\) clipped integers, not a table indexed by all possible sequences of amendments. The bounds can be extracted by a single pass over a supplied core prefix tree; obtaining that tree or constructing a succinct equivalent may be expensive.

If the old tests included multiple separately allowed protocols, intersect their boxes. If a subsequent authorized change admits a new core on the same named quantities, the continuations admitted using \(u\) are exactly those guaranteed by the old proof alone over this box. Those admitted using \(\ell\) are exactly those possible in at least one compatible history. This follows because relaxing expiries can only add words in the new core and both endpoints are members of the box. Restricting the old core while keeping its test map never needs new expiry evidence. After an expansion, a newly enabled test whose slot lies between the retained endpoints requires further evidence to decide it when all other tests on the proposed continuation already pass. The two-endpoint result is for abstract expiry vectors. Intersect with the actual source-history constraints before using it physically; if an endpoint is not realized, it may only give a bound. These are tests for specified continuations, not a policy synthesis theorem.

\paragraph{Indivisible checks.}
If one step checks several old expiries at once, admitted words still impose componentwise upper bounds. But the first refused extension says only that at least one of its tests fails. The exact class is the upper box intersected with one disjunctive constraint per first refusal. For example, at time zero an indivisible grant validates both tags. The histories \((0,1)\) and \((1,0)\) both refuse; their coordinatewise minimum passes and wrongly licenses the grant. With \(n\) separately available grants testing disjoint tag pairs, the class refusing all grants can have \(2^n\) incomparable minimal members. This counts endpoints of the class, not bits: the same class in this construction has just \(n\) two-test clauses. The upper endpoint still exists and gives exact guaranteed continuations over the abstract class. To test possibility of a specified new word, replace each retained upper endpoint by the smaller of itself and the first time the word tests that coordinate. The resulting vector is the largest one compatible with the proposed word. If it violates an old refusal clause, reducing it cannot restore a blocker; if it satisfies all such clauses it is itself a witness. This checks possibility without enumerating the minimal endpoints. The example requires truly indivisible observation of the grant; if the component checks have their own meaningful prefixes, retaining those separate results gives a richer interface.

\paragraph{Source-model test.}
Use the existing \href{https://github.com/eternal-november/eos-research-programmes/pull/390}{two-stream port code} and the actual old mint histories \((-9,-8,-5,-2)\) and \((-9,-7,-5,-2)\) from \href{https://github.com/eternal-november/eos-research-programmes/pull/448}{the fixed-menu study}. The same tag cannot be reused for ten slots. The two full prefix languages agree under each separate fixed timing law even with a competing second stream. Their intersection of boxes admits a second expiry between zero and three before enforcing the actual mint order. The first expiry is exactly one; strictly ordered actual mints tighten the second to two or three.

The new rule allows the second gap law only after both streams have actually made a transfer. Let the partner transfer at slot 0 and the old stream transfer at slot 1. Under the new gap law the old stream can make its second transfer at slot 2 in exactly the earlier-expiry history. The first gap law alone would require slot 3 even with the earlier expiry. The joint core yields 3,574 versus 3,480 admitted labeled prefixes in the two histories. The port source and a separate direct clock calculation agree on the entire bounded languages. Without further information the decision at slot 2 remains ambiguous; a qualified predecision report of \(e_1\le2\) resolves it in this pair (indices start at zero). A partner's actual transfer enables the law; it does not prove the tag is clear. The example implements the authority condition as a test driver around the existing port model, not an independently selected physical workload.

The replay checks 1,772 interval descriptions and post-change word bounds across 256 finite cores, as well as the atomic obstruction and exact selected model. Lean checks the first-refusal characterization, closure under coordinatewise minimum for serial tests, and failure of that closure in the two-test atomic example. An exact direct check with the same history and evidence agrees. This is a reuse theorem under a declared interface and an authored application, not a claim of independent transfer, positive architectural advantage, or full R2 acceptance.
\end{document}
